\appendix %
\section*{Appendix}\section{Results from higher order Fourier analysis}
In this appendix we collect some results from higher order Fourier analysis which we need at various points in the paper. All are known to experts and/or are small perturbations of results that already appear in the literature.

\begin{thm}[Simultaneous arithmetic regularity lemma]\label{t:arl}
Let $f_1, \ldots, f_t: [N] \to [0,1]$ be functions, let $s \geq 1$ be an integer, let $\eps > 0$, and let $\cF: \R^+ \to \R^+$ be a growth function. Then there exists a quantity $M=O_{s,\eps,\cF,t}(1)$ and  decompositions 
\[f_i = f_{i,\nil} + f_{i,\sml} + f_{i,\unf} \]
of each $f_i$ into functions $f_{i, \nil}, f_{i, \sml}, f_{i,\unf}: [N] \to [-1,1]$ where 
\begin{enumerate}
\item($f_{i,\nil}$ structured) $f_{i,\nil}$ is a $(\cF(M), N)$-irrational virtual nilsequence of degree $\leq s$, complexity $\leq M$ and scale $N$, and we may write $f_{i,nil}(n) = F_i(g(n) \Gamma, n \ (\mo q), n/N)$, where $g$ is a polynomial sequence adapted to a filtered nilmanifold and $\Gamma$ is a cocompact lattice on that nilmanifold,
\item($f_{i,\sml}$ small) $f_{i,\sml}$ has an $L^2[N]$-norm of at most $\eps$,
\item($f_{i,\unf}$ very uniform) $f_{i,\unf}$ has a $U^{s+1}[N]$-norm of at most $1/\cF(M)$,
\item(Nonnegativity) $f_{i,\nil}$ and $f_{i,\nil} + f_{i,\sml}$ take values in $[0,1]$.
\end{enumerate}
In particular, the  functions $\{f_{i,\nil}\}_{i=1}^t$ differ only by the choice of Lipschitz function $\{F_i\}_{i=1}^t$. 
\end{thm}
\begin{proof}
This is \cite[Theorem 1.2]{GT10} together with the extra ingredient that the nilsequences $f_{i,\nil}$ differ only by the choice of Lipschitz function. The existence of  this result is well-known to experts. 

One derives the above statement from \cite{GT10} by invoking the non-irrational regularity lemma \cite[Proposition 2.7]{GT10} at each $i$, forming the product polynomial sequence $g(n) = (g_1(n), \ldots, g_t(n))$ on the product nilmanifold $\prod_{i=1}^t G_i/\prod_{i=1}^t \Gamma_i$ and then factorising for irrationality as in \cite[Chapter 2]{GT10}. We omit the details but remark that they are worked out in a similar setting in \cite[Appendix A]{K20}.
\end{proof}

We refer the reader to \cite{GT10} for further background on the theorem that follows. 

\begin{thm}[Flag counting lemma, {\cite[Theorem 1.11]{GT10}}]\label{t:cl}
Let $M, D, t, s$ be positive integers with $D,t,s \leq M$, let $(G/\Gamma, G_\bullet)$ be a degree $\leq s$ filtered nilmanifold of complexity $\leq M$,  let $g: \Z \to G$ be an $(A,N)$-irrational polynomial sequence adapted to $G_\bullet$, let $\Psi = (\psi_1,\ldots, \psi_t)$ be a system of linear forms  with coefficients of magnitude at most $M$ and which satisfies the flag condition, let $P$ be a convex subset of $[-N,N]^D$ with positive Lebesgue measure, let $\Lambda \leq \Z^D$ be a sublattice of index $[\Z^D:\Lambda]\leq M$ and let $\vect x_0 \in \Z^D$. Then for any sequence of $M$-Lipschitz functions $F_1,\ldots,F_t: G/\Gamma \to \C$, we have 
\[\E_{\vect x \in (\vect x_0 + \Lambda )\cap P} \prod_{i=1}^t F_i(g(\psi_i(\vect x))\Gamma) =  \int_{g(0)^\Delta G^\Psi/\Gamma^\Psi}\otimes_{i=1}^t F_i + o_{A\to \infty;M}(1) + o_{N \to \infty;M}(1),\]
where $g(0)^\Delta:= (g(0),\ldots, g(0)) \in G^t$.
\end{thm}

Recall the definition of the \textit{Cauchy-Schwarz complexity} of a system of linear forms (cf. \cite[Definition 1.3.2]{T12}). 

\begin{thm}[Generalised von Neumann inequality on convex sets]\label{t:gvN}
Let $t,D$ be positive integers and let $L>0$. Let $f_1, \ldots, f_t:[N] \to \C$ be 1-bounded. Let $\Psi = (\psi_1, \ldots, \psi_t)$ be a system of affine-linear forms with constant term of size at most $LN$ and linear coefficients of size at most $L$. Let $K \subset [N]^D$ be convex and such that $\Psi(K) \subset [N]^t$. Then letting $s$ be the Cauchy-Schwarz complexity of $\Psi$ and letting $\delta$ be such that $\min_{i=1}^t||f_i||_{U^{s+1}[N]} \leq \delta$ we have
\[ \E_{\vect x \in K} \prod_{i=1}^t f_i(\psi_i(\vect x)) = o_{\delta \to 0;L}(1) + o_{N \to \infty;\delta,L}(1).\] 
\end{thm}
\begin{proof} This is essentially \cite[Proposition 7.1]{GT10L}, except we are dealing with functions $f_i$ which are 1-bounded rather than bounded by a pseudorandom measure. The proof of  \cite[Proposition 7.1]{GT10L} which is given in  \cite[Appendix C]{GT10L} is easily adapted to recover the statement above.
\end{proof}



\begin{prop}[Gowers uniformity on subprogressions]\label{p:gowersUniformSubprog}
Let $f:[N]\to \C$ be 1-bounded, let $\eps >0$ and let $P$ be a subprogression of $[N]$ of size $\eps N$. Then \[||f||_{U^{s+1}(P)} = O_{N \to \infty;\eps}(||f||_{U^{s+1}[N]}) + o_{N \to \infty;\eps}(1).\]
\end{prop}
\begin{proof}
In this proof, all asymptotic notation is implicitly with respect to the limit $N \to \infty$. From Subsection \ref{ss:not-conv} we have $||f||_{U^{s+1}(P)} = \Theta_\eps(||1_Pf||_{U^{s+1}[N]})$. Let $I_\eps = [-\eps N/2,\eps N/2]$. Write $P=qI_\eps+c$. Let $\chi_0$ be the function $\Z \to \R$ which is $1$ on $I_\eps$, which is zero when $|x| > (1+1/\sqrt {N})\eps N/2$, and which linearly interpolates otherwise ($\chi_0$ is the Fourier transform of a de la Vall\'ee Poussin kernel). It is standard that $||\hat \chi_0||_1 = O_\eps(1)$ (one may see this by writing $\chi_0$ as the convolution of suitably normalised characteristic functions of intervals).   
Define $\chi$ by $\chi(x) = \chi_0((x-c)/q)$ for all $x = c\ (\mo q)$ and zero otherwise. Then we have $||\hat \chi||_1 = O_\eps(1)$ and furthermore that $1_P-\chi$ is supported on a set of size $o_\eps(N)$. The former property and Fourier inversion yield that 
\[ ||\chi f||_{U^{s+1}[N]} \leq ||\hat \chi||_1\sup_\theta||e(\theta \cdot)f||_{U^{s+1}[N]} = O_\eps(||f||_{U^{s+1}[N]}),\]
using that the $U^{s+1}[N]$-norm is linear phase invariant (see \cite[Exercise 1.3.21]{T12}). The latter property yields that $||(1_P-\chi)f||_{U^{s+1}[N]} = o_\eps(1)$. We are ready to conclude: 
\[ ||1_Pf||_{U^{s+1}[N]} \leq ||\chi f||_{U^{s+1}[N]} + ||(1_P-\chi)f||_{U^{s+1}[N]} = O_{\eps}(||f||_{U^{s+1}[N]}) + o_{\eps}(1).\]
\end{proof}
